%% SCRAP: experiments/02-experiments/factorial-doe/workflow
%% SOURCE: docs/working/experiments/02-experiments/factorial-doe/workflow.md
%% STATUS: CURRENT
%% FITS: experiments/ch-factorial, cookbook/ch-doe
%% EDITORIAL: lifted — prose rewritten to press voice

\section{Factorial DoE Analysis Workflow}
\label{sec:factorial-doe-workflow}

\subsection{End-to-End Process}

The factorial experiment produces a single CSV file in the StarForth
repository that is then analysed in a separate environment:

\begin{enumerate}
  \item Run \texttt{run\_factorial\_doe.sh} — 64 builds, 1{,}920+ randomised
        runs, CSV and logs written to the experiment label directory.
  \item Transfer the CSV to the analysis environment.
  \item Load into Python or R; separate factor columns from metric columns.
  \item Compute main effects and two-way interactions.
  \item Identify optimal configuration subsets and generate visualisations.
  \item Write the summary report.
\end{enumerate}

\subsection{Phase 1: Data Collection}

The run script handles all 64 configurations:

\begin{lstlisting}[language=bash]
# Quick validation (30 min)
./scripts/run_factorial_doe.sh --runs-per-config 1 TEST_01

# Full production run (2-4 hours)
./scripts/run_factorial_doe.sh --runs-per-config 30 FULL_FACTORIAL
\end{lstlisting}

%% TODO(bob): confirm canonical path for run_factorial_doe.sh in published repo

Monitor progress while the script runs:

\begin{lstlisting}[language=bash]
watch -n 5 'wc -l /path/to/experiment_results.csv'
\end{lstlisting}

Expected progression for a 30-replicate run: approximately 500 rows after
30~minutes, 1{,}500 rows after two hours, 1{,}921 rows at completion.

\subsection{Phase 2: Data Loading}

\begin{lstlisting}[language=Python]
import pandas as pd
import numpy as np

df = pd.read_csv("experiment_results.csv")
print(f"Loaded {len(df)} runs, "
      f"{len(df.groupby('configuration'))} configurations")

factor_cols = ['enable_loop_1_heat_tracking',
               'enable_loop_2_rolling_window',
               'enable_loop_3_linear_decay',
               'enable_loop_4_pipelining_metrics',
               'enable_loop_5_window_inference',
               'enable_loop_6_decay_inference']

metric = 'vm_workload_duration_ns_q48'
\end{lstlisting}

%% TODO(bob): verify column names against actual CSV header from a completed run

\subsection{Phase 3: Main Effects}

For each loop, the main effect is the average metric change when the loop
transitions from off to on, marginalised over all other loop states:

\begin{lstlisting}[language=Python]
main_effects = {}
for loop in factor_cols:
    on_mean  = df[df[loop] == 1][metric].mean()
    off_mean = df[df[loop] == 0][metric].mean()
    effect   = on_mean - off_mean
    pct      = 100.0 * effect / off_mean
    main_effects[loop] = {'effect': effect, 'pct': pct}

for loop, e in sorted(main_effects.items(),
                      key=lambda x: abs(x[1]['pct']), reverse=True):
    print(f"{loop:40} {e['pct']:+7.2f}%")
\end{lstlisting}

\subsection{Phase 4: Two-Way Interactions}

The interaction effect for loops $a$ and $b$:

\begin{equation}
  \hat{\beta}_{ab} = \bar{y}_{(1,1)} - \bar{y}_{(1,0)} - \bar{y}_{(0,1)} + \bar{y}_{(0,0)}
\end{equation}

A positive $\hat{\beta}_{ab}$ indicates synergy; negative indicates suppression.

\begin{lstlisting}[language=Python]
from itertools import combinations

for loop_a, loop_b in combinations(factor_cols, 2):
    y_00 = df[(df[loop_a]==0)&(df[loop_b]==0)][metric].mean()
    y_01 = df[(df[loop_a]==0)&(df[loop_b]==1)][metric].mean()
    y_10 = df[(df[loop_a]==1)&(df[loop_b]==0)][metric].mean()
    y_11 = df[(df[loop_a]==1)&(df[loop_b]==1)][metric].mean()
    interaction = y_11 - y_10 - y_01 + y_00
    if abs(interaction) > 100_000:
        print(f"{loop_a} x {loop_b}: {interaction:+,.0f}")
\end{lstlisting}

\subsection{Phase 5: Optimal Configuration Search}

\begin{lstlisting}[language=Python]
# Top 10 configurations by mean performance
best = (df.groupby('configuration')[metric]
          .mean()
          .nlargest(10)
          .reset_index())
print(best.to_string())
\end{lstlisting}

\subsection{Phase 6: Reporting}

The analysis report covers four sections:

\begin{enumerate}
  \item \textbf{Main effects table} — individual loop impacts sorted by effect
        size with $p$-values.
  \item \textbf{Interaction matrix} — $6 \times 6$ table of two-way interaction
        estimates; cells highlighted for synergy or suppression.
  \item \textbf{Optimal configurations} — top 3--5 production candidates on the
        Pareto frontier of performance versus code complexity.
  \item \textbf{Statistical summary} — variance by configuration, 95\%
        confidence intervals, regression $R^2$.
\end{enumerate}

